\documentclass[titlepage]{ctexart}

\usepackage{amsmath}
\usepackage{amssymb}
\usepackage{enumerate}

\renewcommand\parallel{\mathrel{/\mskip-2.5mu/}}

\title{第六章~立体几何初步}
\author{dovego}
\date{\today}

\begin{document}
\maketitle


\section{基本立体图形}

\subsection{构成空间几何体的基本元素}

(顶)点, 线(棱), 面.

\subsection{简单多面体——棱柱、棱锥和棱台}

多面体: 由平面多边形围成的几何体. 多这些边形称为多面体的面, 两个相邻面的公共边称为多面体的棱, 棱与棱的公共点称为多面体的顶点.

\subsubsection{棱柱}

棱柱: 有两个面边数相同且互相平行; 其余各面都是四边形, 且每相邻两个面的公共边都互相平行. 

在棱柱中, 两个互相平行的面称为\textbf{底面}, 简称\textbf{底}; 其余各面称为\textbf{侧面}; 相邻侧面的公共边称为\textbf{侧棱};
侧面与底面的公共顶点称为\textbf{顶点}; 既不在同一底面也不在同一侧面的两个顶点的连线称为\textbf{对角线}; 
过上底面一点 $O_1$ 作下底面的垂线, 垂足为 $O$, 则距离 $OO_1$ 称为棱柱的\textbf{高}.

棱柱既可以用两个底面各顶点表示为\textbf{棱柱 $ABCDE-A_1B_1C_1D_1E_1$}, 也可以用某条对角线的两个端点表示为\textbf{棱柱 $AC_1$}.

棱柱的性质: 
\begin{enumerate}[(1)]
    \item 侧棱都相等;
    \item 两个底面和平行于底面的截面都是全等的多边形;
    \item 过不相邻两条侧棱的截面都是平行四边形.
\end{enumerate}

侧面都是矩形的棱柱称为\textbf{直棱柱}, 其他的称为\textbf{斜棱柱}. 底面是正多边形的直棱柱称为\textbf{正棱柱}.

底面是平行四边形的棱柱称为\textbf{平行六面体}.

\subsubsection{棱锥和棱台}

一个面是多边形, 其余各面都是有一个公共顶点的三角形围成的几何体称为\textbf{棱锥}.

底面是正多边形, 且顶点在过底面中心且与底面垂直的直线上的棱锥称为\textbf{正棱锥}. 
它的侧面都是全等的等腰三角形, 这些等腰三角形底边上的高称为正棱锥的\textbf{斜高}.

(棱锥的性质) 平行于底面的截面与底面相似.

用平行于底面的平面去截棱锥, 截面与底面之间的部分称为\textbf{棱台}.
棱台有上底面, 下底面, 侧面, 侧棱, 高. 

由正棱锥截得的棱台称为\textbf{正棱台}. 它的侧面都是全等的等腰梯形, 这些等腰梯形的高称为正棱台的\textbf{斜高}.

\subsection{简单旋转体——球、圆柱、圆锥和圆台}

一条平面曲线绕其所在平面内一条定直线旋转一周所得曲面称为\textbf{旋转面}, 封闭的旋转面围成的几何体称为\textbf{旋转体}.

\textbf{球}: 球面, 球体, 球心, 半径, 直径.

圆柱, 圆锥, 圆台: 底面, 侧面, 高, 旋转轴, 母线. 它们过旋转轴的截面都是全等的(等腰)图形.


\section{直观图}

一些常用画法. 平行投影法: 斜二测(多面体), 正等测(圆), 中心投影法.

\textbf{斜二测画法}的原理: 在新的坐标系中, 使 $\angle x'Oy' = 45^\circ$ 保持各点的横坐标不变, 将纵坐标变为原来的一半.

几何体的斜二测画法原理: 在新的坐标系中, 使 $\angle x'Oy' = 45^\circ, \angle xO'z = \angle y'Oz = 90^circ$,
保持各点 $x', z'$ 坐标和原来一样, $y'$ 坐标变为原来的一半. 


\section{空间点、直线、平面之间的位置关系}

\subsection{空间图形基本位置关系的认识}

\subsubsection{点与直线、点与平面的位置关系}

点在直线上: $B \in b$. 点在直线外: $B \notin b$.

点在平面内: $B \in \alpha$. 点在平面外: $B \notin \alpha$.

\subsubsection{直线与直线、直线与平面的位置关系}

直线与直线相交: $a \cap l = B$. 不相交: $b \cap l = \varnothing$.

直线在平面内: $a \subset \alpha$. 与平面相交: $a \cap \alpha = A$. 
与平面平行: $a \parallel \alpha \Leftrightarrow a \cap \alpha = \varnothing$.

\subsubsection{平面与平面的位置关系}

平面与平面不相交(平行): $\alpha \parallel \beta \Leftrightarrow \alpha \cap \beta = \varnothing$.
相交: $\alpha \cap \beta = l$.

\subsection{刻画空间点、线、面位置关系的公理}

\textbf{基本事实 1} \quad 过不在一条直线上的三个点, 有且只有一个平面.

\textbf{基本事实 2} \quad 如果一条直线上的两个点在一个平面内, 那么这条直线在这个平面内.

\textbf{推论 1} \quad 一条直线和该直线外一点确定一个平面.

\textbf{推论 2} \quad 两条相交直线确定一个平面.

\textbf{推论 3} \quad 两条平行直线确定一个平面.

\textbf{基本事实 3} \quad 若两个不重合的平面有一个公共点, 那么它们有且仅有一条过该点的公共直线.
即 $P \in \alpha, P \in \beta \Leftrightarrow \alpha \cap \beta = l, \text{且}P \in l$.

\textbf{基本事实 4} \quad 平行于同一直线的两条直线互相平行.

不同在任何一个平面内的两条直线称为\textbf{异面直线}.

\textbf{定理} \quad 如果空间中两个角的两条边分别对应平行, 那么这两个角相等或互补.

已知两条异面直线 $a, b$, 过空间中的任意一点 $O$ 作直线 $a' \parallel a, b' \parallel b$,
直线 $a', b'$ 所成的不大于 $90^\circ$ 的角称为\textbf{异面直线 $a, b$ 的夹角}.
若夹角为 $90^\circ$, 则称异面直线 $a \perp b$.

四个顶点不在同一平面的四边形称为\textbf{空间四边形}.


\section{平行关系}

\subsection{直线与平面平行}

\subsubsection{直线与平面平行的性质}

\textbf{性质定理} \quad 一条直线与一个平面平行, 若过该直线的平面与此平面相交, 则该直线与交线平行.

\subsubsection{直线与平面平行的判定}

\textbf{判定定理} \quad 若平面外一条直线与此平面内一条直线平行, 那么该直线与此平面平行.

\subsection{平面与平面平行}

\subsubsection{平面和平面平行的性质}

\textbf{性质定理} \quad 两个平面平面, 若另一个平面与这两个平面相交, 那么两条交线平行.

\subsubsection{平面和平面平行的判定}

\textbf{判定定理} \quad 如果一个平面内的两条相交直线与另一个平面平行, 那么这两个平面平行.
即 $a \subset \alpha, b \subset \alpha, a \cap b = A, a \parallel \beta, b \parallel \beta \Rightarrow \alpha \parallel \beta$.


\section{垂直关系}

\subsection{直线与平面垂直}

一般地, 如果一条直线与平面内的所有直线都垂直, 则称这条直线与这个平面垂直, 记作 $l \perp \alpha$.
直线 $l$ 称为平面 $\alpha$ 的\textbf{垂线}, 平面 $\alpha$ 称为直线 $l$ 的\textbf{垂面}, 它们唯一的公共点 $P$ 称为\textbf{垂足}.

过一点有且仅有一条直线与一个平面垂直, 过一点有且仅有一个平面与一条直线垂直.

\subsubsection{直线与平面垂直的性质}

\textbf{性质定理} \quad 垂直于同一平面的两条直线平行.

如果一条直线于平面平行, 那么这条直线上的任意一点到平面的距离就是这条\textbf{直线到这个平面的距离}.

平面的一条\textbf{斜线}与它在平面上的投影所成的\textbf{锐角}, 叫做\textbf{直线与这个平面的夹角}.
若直线垂直于平面, 则称夹角为直角. 若直线平行于平面或在平面内, 则称夹角为 $0^\circ$.

\subsubsection{直线和平面垂直的判定}

\textbf{判定定理} \quad 若一条直线与平面内的两条相交直线垂直, 则这条直线与该平面垂直.

\subsection{平面与平面垂直}

半平面, 二面角, 二面角的棱, 二面角的面, 二面角的平面角, 直二面角.

\subsubsection{平面与平面垂直的性质}

两个平面垂直, 若一个平面有一条直线垂直于这两个平面的交线, 则这条直线与另一个平面垂直.

\subsubsection{平面与平面垂直的判定}

\textbf{判定定理} \quad 如果一个平面过另一个平面的垂线, 那么这两个平面垂直.
即 $l \subset \alpha, l \perp \beta \Rightarrow \alpha \perp \beta$.


\section{简单几何体的再认识}

\subsection{柱、锥、台的侧面展开与面积}

\subsubsection{圆柱、圆锥与圆台}

侧面积的公式: 
\begin{align*}
    S_{\text{圆柱侧}} = 2\pi rl, && S_{\text{圆锥侧}} = \pi rl, && S_{\text{圆台侧}} = \pi(r_1 + r_2)l.
\end{align*}

\subsubsection{直棱柱、正棱锥、正棱台}

侧面积的公式($c$ 为底面周长, $h'$ 为斜高):
\begin{align*}
    S_{\text{直棱柱侧}} = ch, && S_{\text{正棱锥侧}} = \dfrac{1}{2}ch', && S_{\text{正棱台侧}} = \dfrac{1}{2}(c_1 + c_2)h'.
\end{align*}

\subsection{柱、锥、台的体积}

棱柱和圆柱: $V_{\text{柱体}} = Sh$.
$S$ 为柱体底面积, $h$ 为柱体的高.

棱锥和圆锥: $V_{\text{锥体}} = \dfrac{1}{3}Sh$.
$S$ 为锥体底面积, $h$ 为锥体的高.

棱台和圆台: $V_{\text{台体}} = \dfrac{1}{3}(S_{\text{上}} + S_{\text{下}} + \sqrt{S_{\text{上}} \cdot S_{\text{下}}})h$.
$h$ 为台体的高.

\subsection{球的表面积和体积}

过球心的平面截得的圆称为\textbf{球的大圆}, 不过球心的称为\textbf{球的小圆}.

当直线与球有唯一交点时, 称直线与圆相切, 称这一点为切点.

过球外一点的所有切线的切线长都相等.

球的表面积和体积公式
\begin{align*}
    S_{\text{球面}} = 4\pi R^2, && V_{\text{球体}} = \dfrac{4}{3}\pi R^3.
\end{align*}


\end{document}